\section{问题三：全向干扰源的自动搜索、定位与清除}\label{sec:q3}

\subsection{三层结构与终止证明}

第三问的难点不在“怎样清除一个已知的源”，而在“怎样证明已经没有源了”。本文把策略分成\textbf{发现—定位—清除}三层，其中发现层承担终止证明的责任。

\begin{definition}[距离覆盖测站集]\label{def:cover3}
称测站集 $\mathcal S$ 对目标区域 $\Dreg$ 具有\emph{距离覆盖}，若
\begin{equation}\label{eq:cover3}
  d_{\max}(\mathcal S):=\max_{x\in\Dreg}\ \min_{s\in\mathcal S}\ \|x-s\|\ \le\ \rho_{\min}=1000\mmet .
\end{equation}
\end{definition}

\begin{theorem}[终止的正确性]\label{thm:term3}
设 $\mathcal S$ 具有距离覆盖。若算法在每个 $s\in\mathcal S$ 处都对每个尚未成功清除的频道检测过一次，并清除了由此发现的所有目标，则 $\Dreg$ 内不存在未被清除的全向干扰源。
\end{theorem}
\begin{proof}
反设频道 $c$ 上还存在未清除源 $g\in\Dreg$。由\cref{eq:cover3}存在 $s\in\mathcal S$ 使 $\|s-g\|\le1000\le\rho$。全向源覆盖角为 $360^\circ$，故在 $s$ 处检测频道 $c$ 必定返回 \texttt{direction} 或 \texttt{near}，于是 $g$ 会被发现并进入目标列表，与“所有发现的目标均已清除”矛盾。
\end{proof}

\cref{thm:term3}说明：合法的停止理由只有两个——\textbf{(i)} 覆盖测站全部访问完毕且目标列表清空；\textbf{(ii)} 已成功清除 $16$ 个源，达到题设数量上界。“连续若干次没有收到信号”“已经清了 $10$ 个”“$20$ 个频道各扫过一遍”都\emph{不是}有效的终止证明。

\subsection{发现层：七站布局及其最优性}\label{sec:q3:layout}

\subsubsection{最坏覆盖距离的解析式}

由\cref{sec:analysis:fixed}，固定成本与测站数成正比，因此应在\cref{eq:cover3}的约束下尽量减少站数、缩短巡回。本文采用“圆心 $+$ 半径 $r$ 的正六边形”共 $7$ 站的布局
\begin{equation}\label{eq:hex}
  \mathcal S_7=\{\bm 0\}\cup\bigl\{r\,\uvec{60^\circ k}\ :\ k=0,1,\dots,5\bigr\}.
\end{equation}

\begin{theorem}[七站布局的最坏覆盖距离]\label{thm:hexcover}
对 $0<r\le\sqrt3R$，
\begin{equation}\label{eq:dmax}
  d_{\max}(\mathcal S_7)=\max\left\{\frac{r}{\sqrt3},\ \sqrt{R^2+r^2-\sqrt3\,Rr}\right\}.
\end{equation}
\end{theorem}
\begin{proof}
设 $x\in\Dreg$，$t=\|x\|$。六个环点均匀分布，故存在一个环点 $s_k$ 使 $x$ 与 $s_k$ 的极角差 $\phi$ 满足 $|\phi|\le30^\circ$，于是
$\|x-s_k\|=\sqrt{t^2+r^2-2tr\cos\phi}\le\sqrt{t^2+r^2-\sqrt3\,tr}=:g(t)$。
又圆心站给出距离 $t$，故最近站距离不超过 $f(t)=\min\{t,g(t)\}$。

比较两支：$t\le g(t)\iff t^2\le t^2+r^2-\sqrt3tr\iff t\le r/\sqrt3$。因此
$f(t)=t$ 于 $[0,r/\sqrt3]$，$f(t)=g(t)$ 于 $[r/\sqrt3,R]$。前段最大值为 $r/\sqrt3$。后段 $g$ 是关于 $t$ 的二次根式，$g'(t)\propto2t-\sqrt3r$，在 $t=\sqrt3r/2$ 取极小，故在闭区间 $[r/\sqrt3,R]$ 上的最大值在端点取到，为 $\max\{g(r/\sqrt3),g(R)\}=\max\{r/\sqrt3,\ \sqrt{R^2+r^2-\sqrt3Rr}\}$。合并即得\cref{eq:dmax}。上界可达：$x$ 取 $t=r/\sqrt3$ 且极角恰在两环点之间（Voronoi 三重点），或 $x$ 取目标圆边界上两环点角平分方向的点。
\end{proof}

\begin{corollary}[环半径的可行区间与最优值]\label{cor:ring}
取 $R=1800\mmet$、$\rho_{\min}=1000\mmet$，$d_{\max}\le\rho_{\min}$ 当且仅当
\begin{equation}\label{eq:ringrange}
  \begin{aligned}
  &r\in\bigl[\,r_-,\ r_+\,\bigr],\\
  &r_-=\frac{\sqrt3R-\sqrt{3R^2-4(R^2-\rho_{\min}^2)}}{2}=1122.956\mmet,\\
  &r_+=\sqrt3\,\rho_{\min}=1732.051\mmet .
  \end{aligned}
\end{equation}
$d_{\max}(r)$ 在 $r^\star=\dfrac{\sqrt3}{2}R=1558.846\mmet$ 处取最小值 $R/2=900\mmet$。
\end{corollary}

\paragraph{为什么不取覆盖最优的 $r^\star$.} 覆盖只是\emph{约束}而非\emph{目标}。$r$ 越小，六个环点越靠内，巡站路线越短，而且发现目标后的第二测点、补测点与清除点也都更靠近主路线。因此本文在可行区间内取靠近下界的 $r=1150\mmet$，此时
\begin{equation}
  d_{\max}=\sqrt{1800^2+1150^2-\sqrt3\times1800\times1150}=988.5114\mmet<1000\mmet,
\end{equation}
覆盖余量 $11.4886\mmet$。$5\mmet$ 网格加边界采样的数值核对（$61445$ 个位置）与该解析值一致。程序在启动时用\cref{eq:dmax}重新校验，不满足覆盖条件即\emph{拒绝运行}，避免误改参数导致保证失效。

\cref{fig:q3-layout}给出了该布局的覆盖示意与 $d_{\max}(r)$ 曲线。实验上，$r$ 取 $1250,1300,\dots,1450$ 的扫描全部慢于 $1150$；而 $r<1122.956$ 会直接破坏覆盖保证。也就是说，\textbf{$r=1150\mmet$ 已经贴着可行下界}，这一维参数上不再有可挖掘的空间。

\subsubsection{七站是保证发现的最小站数}

\begin{theorem}[站数下界]\label{thm:seven}
在 $R=1800\mmet$、$\rho_{\min}=1000\mmet$ 的条件下，任何具有距离覆盖的测站集至少包含 $7$ 个站；\cref{eq:hex}取 $r\in[r_-,r_+]$ 时达到该下界。
\end{theorem}
\begin{proof}
\cref{eq:cover3}等价于用 $|\mathcal S|$ 个半径 $\rho_{\min}$ 的闭圆盘覆盖半径 $R$ 的圆盘，即覆盖半径比为 $R/\rho_{\min}=1.8$。Bezdek 证明了 $6$ 个单位圆盘所能覆盖的圆的最大半径为
$1.7988\ldots<1.8$\cite{bezdek,bezdek2,kershner}，故 $|\mathcal S|\le6$ 不可能；而 $\mathcal S_7$ 在 $r\in[r_-,r_+]$ 时满足覆盖，故 $7$ 可达。
\end{proof}

为了从反面确认这一结论的紧性（并检验“少一站能省多少”），本文另外实现了一个\emph{不保证覆盖}的六站实验入口，结果见\cref{tab:q3six}。

\begin{table}[!htbp]
  \caption{六站布局实验（独立入口，默认方案未采用）}\label{tab:q3six}\centering\small
  \begin{tabularx}{\textwidth}{@{}L c c c c c@{}}
    \toprule[1.5pt]
    布局 & 最坏未覆盖距离/m & 单源漏检概率 & $120$ 例均值 s/源 & 与 $7$ 站配对差 & 更快局数 \\
    \midrule[1pt]
    $7$ 站（默认） & $988.5$（有保证） & $0$ & $258.6$ & --- & --- \\
    $6$ 站环 $960\mmet$   & $1102.0$ & $\approx4\times10^{-4}$ & $256.1$ & $-2.5$ & $69/120$（漏 $1$ 源） \\
    $6$ 站环 $1000\mmet$  & $1062.9$ & $\approx2\times10^{-4}$ & $257.2$ & $-1.4$ & $61/120$（无漏） \\
    $6$ 站环 $1040\mmet$  & $1043.5$ & $\approx8\times10^{-5}$ & $257.9$ & $-0.7$ & $63/120$（无漏） \\
    \bottomrule[1.5pt]
  \end{tabularx}
\end{table}

用 Nelder--Mead 做过六站极小极大优化，最坏未覆盖距离最低只能压到约 $1003\mmet$（接近理论值 $1000.7\mmet$），仍然超过 $1000\mmet$。更重要的是：\textbf{去掉一个站只节省约 $6\msec/$源的扫描，却因为失去圆心站而使首次发现更远、清除路程更长，移动增加约 $4\msec/$源，净收益仅 $1\sim2\msec/$源，代价是引入约万分之二的单源漏检概率}。在“清除率是硬约束”的准则下这一交换不可接受，因此正式方案坚持七站。这表明\textbf{第三问的发现层已接近其几何下界}。

\begin{figure}[!htbp]
  \centering
  \fbox{\parbox[c][6.5cm][c]{0.9\textwidth}{\centering\textbf{【此处插图】图~\thefigure：第三问七站布局与覆盖验证}\\[4pt]
  \small 画法建议：左子图画 $R=1800\mmet$ 目标圆、$7$ 个测站（圆心 $+$ 环半径 $1150\mmet$ 的六点）与它们的 Voronoi 剖分，
  每站叠加半径 $1000\mmet$ 的接收圆（半透明），可见并集完整覆盖目标圆；标出两个最坏点（$r/\sqrt3=663.95\mmet$ 的 Voronoi 三重点
  与边界上 $988.51\mmet$ 的点）。右子图画 $d_{\max}(r)$ 随环半径 $r$ 的变化曲线（\cref{eq:dmax}），
  横轴 $r\in[900,1800]$，标出可行区间 $[1122.96,1732.05]$、最优点 $r^\star=1558.85$、以及实际取值 $r=1150$ 及其余量 $11.49\mmet$。}}
  \caption{七站布局的覆盖证明与环半径的可行区间}\label{fig:q3-layout}
\end{figure}

\subsection{定位层：有界误差可行域}\label{sec:q3:loc}

对每个已发现频道 $c$ 维护一个凸可行域 $P_c$，其不变量是
\begin{equation}\label{eq:invariant}
  \textbf{（包含性不变量）}\qquad g_c\in P_c\ \text{恒成立}.
\end{equation}

\paragraph{（i）初始化.} 首次收到 \texttt{direction} 时，以检测点 $p$ 为顶点构造包含真实扇区的三角形外包
\begin{equation}
  P_c^{(0)}=\conv\bigl\{p,\ p+\rho_{\max}\uvec{\theta}-w\bm n,\ p+\rho_{\max}\uvec{\theta}+w\bm n\bigr\},
  \quad w=\rho_{\max}\tan\bar\varepsilon,\ \bm n\perp\uvec{\theta}.
\end{equation}
它外包了半径 $\rho_{\max}$、半角 $\bar\varepsilon$ 的扇形，因此无需猜测源距。

\paragraph{（ii）新方位裁剪.} 每获得一条新的有效示向度，按\cref{eq:wedge}生成两条半平面并求交：$P_c\leftarrow P_c\cap W$。若结果为空，说明模型假设与观测矛盾，程序\emph{报错并保留日志}，而不是取一个近似点继续。

\paragraph{（iii）目标圆外包约束.} 所有源都在 $\Dreg$ 内，故可以再交上目标圆。为保持多边形表示，本文用 $48$ 条\textbf{支撑半平面}
$\bm n_j^{\!\top}x\le R+10^{-7}$（$\bm n_j$ 均匀取 $48$ 个方向）作为圆的\emph{外切}近似。这里必须是外切而不是内接：内接多边形会把靠近圆边界的一圈真实位置切掉，可能使边界源脱离 $P_c$，破坏\cref{eq:invariant}。

\paragraph{（iv）无信号的距离排除（仅第三问）.} 这是第三问相对第四问独有的信息优势。

\begin{theorem}[无信号排除]\label{thm:nosignal}
第三问中所有源都是全向源。若在点 $s$ 检测频道 $c$ 返回 \texttt{no\_signal}，而频道 $c$ 上仍存在未清除的源 $g_c$，则
$\|s-g_c\|>\rho_c\ge\rho_{\min}$，即 $g_c\notin B(s,\rho_{\min})$。
\end{theorem}

于是可令 $P_c\leftarrow \conv\bigl(P_c\setminus B^\circ(s,\rho_{\min})\bigr)$。注意 $P_c\setminus B^\circ$ 一般\emph{非凸}，本文取其凸包作为\textbf{保守外近似}：这可能重新包含一小部分已被排除的区域（约束变松），但绝不会误删真源。实现中排除半径取 $\rho_{\min}-10^{-4}$，为浮点计算留余量。

\paragraph{一次必须记录的失败.} 在圆盘排除约束的首轮验证中，曾有一例的几何证书失效：该局虽然最终清除了全部源，程序内部却报出“可行域排除了真源”。根因是：重复施加圆盘排除时，一个理论上位于边界的交点被浮点判为落在圆盘\emph{内部}，相应二次方程的根又略微超出线段参数范围，导致边界顶点被错误丢弃，多次裁剪后可行域出现了\emph{不保守}的收缩。修复措施为：用平方距离容差保留边界顶点、对接近端点的根做范围容许与夹取、点数不足或凸包退化时保留输入多边形，并增加重复约束的回归测试。该案例的种子、报错与修复前后对比见附录~\ref{app:debug}。

这次失败给出一条重要教训，也是本文的一个方法论主张：\textbf{“结果全部清除”不能替代对中间几何不变量的检查}。所有实验都必须逐条核验可行多边形与包围圆是否真的包含真源。\cref{fig:q3-shrink}展示了可行域在上述四类约束下逐步收缩直至可以认证清除的完整过程。

\begin{figure}[!htbp]
  \centering
  \fbox{\parbox[c][6.5cm][c]{0.9\textwidth}{\centering\textbf{【此处插图】图~\thefigure：可行域的逐步收缩与认证清除}\\[4pt]
  \small 画法建议：四联子图，同一坐标范围。(a) 首次测向后的三角形外包 $P^{(0)}$（细长扇形外包，长 $1500\mmet$）；
  (b) 第二次测向后 $P^{(0)}\cap W_2$（细长四边形），标出最小包围圆及其半径；
  (c) 叠加一次 \texttt{no\_signal} 的圆盘排除（画出 $B(s,1000)$ 的弧与取凸包后的保守外近似，用两种填充区分“精确差集”与“凸包外近似”，
  说明外近似只会变松、不会误删真源）；(d) 圆心补测后 $\mecr\le20\mmet$，画出认证清除圆与真源位置。
  每个子图标注 $\mecr$ 的数值，并始终把真源标成一个醒目的点以显示包含性不变量始终成立。}}
  \caption{有界误差可行域的收缩过程}\label{fig:q3-shrink}
\end{figure}

\subsection{清除层：认证清除、光学兜底与补测精化}\label{sec:q3:refine}

\paragraph{（i）认证清除.} 由\cref{cor:certify}，当 $\mecr(P_c)\le r_c-10^{-6}$ 时，在最小包围圆圆心 $c_*(P_c)$ 处执行一次 \texttt{/clear} 必定成功。最小包围圆按“$1$、$2$ 或 $3$ 个支撑顶点”穷举候选圆心并逐一校验，凸多边形只需检查顶点。

\paragraph{（ii）光学覆盖兜底.} 若 $\mecr>r_c$，用边长 $h_{\rm opt}=25\mmet$ 的方格覆盖 $P_c$，保留所有与 $P_c$ 相交的格，在各格中心依次尝试清除。覆盖正确性来自
\begin{equation}\label{eq:grid}
  \frac{h_{\rm opt}}{\sqrt2}=\frac{25}{\sqrt2}=17.677670\mmet<20\mmet=r_c,
\end{equation}
即格内任意点到格心的距离都在清除半径内；由于真源必在 $P_c$ 内、$P_c$ 被这些格覆盖，故必有一格成功。遍历采用按列蛇形路径，并根据当前位置选择从哪一端进入。

\paragraph{（iii）补测精化.} 认证清除与光学兜底之间存在一个过渡区间：$\mecr$ 只比 $20\mmet$ 略大（例如 $60\mmet$），却可能需要在上百米的网格上逐点尝试清除。为此加入圆心补测：当 $r_c<\mecr\le300\mmet$ 时，\textbf{先走到 $c_*(P_c)$ 再测一次}——该移动与后续清除路径高度重合，因此边际成本较低；若仍 $>20\mmet$，则沿新方位的垂线方向横移 $60\mmet$ 再补测一次；此后才进入光学网格。$100$ 例留出验证：平均 $262.6\to260.0\msec/$源，\textbf{每源失败清除次数由 $3.00$ 次降到 $0.38$ 次}，$64$ 例更快、$36$ 例更慢、最大单例退步 $30.9\msec/$源。

\subsection{调度层：搜索与目标处理的混合贪心}\label{sec:q3:shared}

早期版本采用“到一个站 $\to$ 扫完 $\to$ 处理本站发现的全部目标 $\to$ 再去下一站”的分批流程，会产生大量折返。当前版本把\textbf{搜索任务}与\textbf{目标任务}放进同一个任务池，每一步比较：

\begin{itemize}
  \item 去最近未访问测站的代价：$\|s^*-p_{\rm cur}\|$；
  \item 处理待清除目标 $c$ 的代价：
  \begin{equation}\label{eq:q3cost}
    J(c)=
    \begin{cases}
      \|q_c-p_{\rm cur}\|+\|q_c-c_*(P_c)\|, & \text{还需补一条方位},\\[2pt]
      \|x_c-p_{\rm cur}\|, & \text{已可直接处理},
    \end{cases}
  \end{equation}
\end{itemize}
其中 $q_c$ 为第二测点，$x_c$ 为 \texttt{near} 点或最小包围圆圆心 $c_*(P_c)$；取二者中较小者执行。这样机器狗可以“顺路”经过新的测站，也可以在靠近时提前把近处目标处理掉。目标状态\emph{跨测站保留}，不再每到一站就清空。

\paragraph{共享测站（shared scan）.} 一次移动应当尽量服务多个目标。到达第二测点或完成一次清除后，算法检查其他待处理频道：若当前位置对该频道\emph{确有观测价值}，就顺便补测一次。筛选条件见\cref{tab:shared}，其作用是防止共享测量无限增加——它们是\emph{选择动作的启发式}，不提供额外的接收保证，也不参与任何正确性论证。

\begin{table}[!htbp]
  \caption{共享测站的筛选条件}\label{tab:shared}\centering
  \begin{tabular}{lc}
    \toprule[1.5pt]
    条件 & 阈值 \\
    \midrule[1pt]
    与该目标已有采样点的最小距离 & $\ge100\mmet$ \\
    目标估计中心到当前位置的距离 & $\le1300\mmet$ \\
    估计交会角的正弦 & $\ge0.12$ \\
    每个目标的额外共享检测次数 & $\le4$ \\
    已有有效方位数 & $<3$ 时才继续 \\
    已清除或已获得 \texttt{near} 的目标 & 跳过 \\
    \bottomrule[1.5pt]
  \end{tabular}
\end{table}

\subsection{算法总览}

把上述四层合在一起，得到\cref{alg:q3}。

\begin{algorithm}[!htbp]
\caption{第三问：七站覆盖 $+$ 混合贪心的在线搜索、定位与清除}\label{alg:q3}
\begin{algorithmic}[1]
\Require 目标圆半径 $R$，接收半径下界 $\rho_{\min}$，误差界 $\bar\varepsilon$，环半径 $r=1150\mmet$
\Ensure 全部干扰源被清除
\State 调用 \texttt{/enter}；$p_{\rm cur}\gets(0,0)$，当前频道 $\gets1$，$\mathcal C_{\rm done}\gets\varnothing$，轨迹表 $\mathcal Q\gets\varnothing$
\State 由\cref{eq:hex}生成 $\mathcal S_7$ 并用\cref{eq:dmax}校验覆盖；不满足则拒绝运行
\While{$\mathcal S_7\ne\varnothing$ 或 $\mathcal Q$ 中仍有未清除频道}
  \State $c^\star\gets\arg\min_{c\in\mathcal Q\setminus\mathcal C_{\rm done}}J(c)$ \Comment{\cref{eq:q3cost}}
  \State $s^\star\gets\arg\min_{s\in\mathcal S_7}\|s-p_{\rm cur}\|$
  \If{$\|s^\star-p_{\rm cur}\|<J(c^\star)$} \Comment{搜索更划算}
    \State 走到 $s^\star$ 并从 $\mathcal S_7$ 移除；对每个未清除频道 $c$ 依次检测（当前频道优先，其余升序）
    \State 若 $c\in\mathcal Q$ 则更新 $P_c$，否则在收到信号时新建轨迹并施加\cref{thm:nosignal}与目标圆约束
  \Else \Comment{处理目标更划算}
    \If{$c^\star$ 尚需补一条方位}
      \State 走到第二测点 $q_{c^\star}=\text{origin}+200\,\uvec{\theta}\pm100\,\bm n$（取综合代价小的一侧）并检测，更新 $P_{c^\star}$
      \State 共享测站扫描（排除 $c^\star$）
    \EndIf
    \State 圆心补测精化（\cref{sec:q3:refine}(iii)）
    \If{$\mecr(P_{c^\star})\le r_c-10^{-6}$} 在 $c_*(P_{c^\star})$ 处清除
    \Else\ 按 $25\mmet$ 方格蛇形遍历 $P_{c^\star}$ 逐格尝试清除
    \EndIf
    \State 共享测站扫描
  \EndIf
  \If{$|\mathcal C_{\rm done}|=16$} \textbf{break} \Comment{题设数量上界}
  \EndIf
\EndWhile
\State 调用 \texttt{/exit}
\end{algorithmic}
\end{algorithm}

\paragraph{虚拟时间上界.} 所有行动点都在以原点为心、$2000\mmet$ 为半径的圆内，任意两点距离不超过 $4000\mmet$。$7$ 个测站每站至多 $20$ 次含切频的检测（$120\msec$）；至多 $16$ 个目标，每个目标至多两段全局跳转、一条保守长度 $6500\mmet$ 的光学路线、若干次检测与至多 $183$ 次清除尝试。于是
\begin{equation}
  T\le 7\Bigl(\frac{4000}{5}+120\Bigr)+16\Bigl[\frac{2\times4000+6500}{5}+6\times6+183\times3+5\Bigr]
  \approx6.2\times10^{4}\msec\ \ll\ 3.6\times10^{5}\msec,
\end{equation}
远小于虚拟时间上限，也说明算法不可能因死循环而被强制终止。实际演练中整局虚拟时间约 $3300\msec$，现实程序运行时间 $1.7\sim2.4\msec$，相对 $1200\msec$ 的限制有三个数量级的余量。

\subsection{实验结果与消融}

\cref{fig:q3-track}给出一局完整任务的行走轨迹，可直观看到巡站与目标处理如何交错进行。所有本地实验均使用与开发阶段\emph{不重叠}的新种子做留出验证，并逐例核验：清除数等于真实源数、每次成功清除的实际距离 $\le20\mmet$、可行多边形与包围圆均包含真源、逐动作时间与\cref{eq:objective}一致。

\begin{table}[!htbp]
  \caption{第三问历次结构性改进的配对实验（单位：虚拟秒/源，两版均 $100\%$ 清除）}\label{tab:q3paired}\centering
  \begin{tabular}{lcccc}
    \toprule[1.5pt]
    改进 & 数据集 & 对照 & 候选 & 下降 \\
    \midrule[1pt]
    单次测向 $\to$ 两次测向 & $40$ 例，$498$ 源 & $977.71$ & $764.71$ & $21.79\%$ \\
    动态处理顺序 $+$ 共享测站 & $100$ 例，$1310$ 源 & $723.82$ & $653.11$ & $9.77\%$ \\
    $25$ 站网格 $\to$ $7$ 站解析覆盖 & $100$ 例，$1308$ 源 & $654.53$ & $346.62$ & $47.04\%$ \\
    环半径 $1150$ $+$ 混合贪心调度 & $100$ 例，$1315$ 源 & $343.73$ & $263.86$ & $23.24\%$ \\
    （同上，压力集） & $70$ 例，$890$ 源 & $336.31$ & $281.83$ & $16.20\%$ \\
    圆心补测精化 & $100$ 例留出 & $262.6$ & $260.0$ & $1.0\%$ \\
    \bottomrule[1.5pt]
  \end{tabular}\\[2pt]
  \footnotesize 各行使用的案例集合不同，\textbf{不能把各行降幅连乘}得到一个“总提升”。
\end{table}

\begin{figure}[!htbp]
  \centering
  \fbox{\parbox[c][6.5cm][c]{0.9\textwidth}{\centering\textbf{【此处插图】图~\thefigure：第三问一局完整轨迹}\\[4pt]
  \small 画法建议：取一局具有代表性的演练（例如 $n=14$、$\bar T=266.6\msec/$源），画出 $1800\mmet$ 目标圆、$7$ 个测站、
  机器狗的完整行走轨迹（按时间用渐变色，起点 $(0,0)$ 标出），每个源的真实位置用叉号、成功清除点用圆圈，
  失败的光学尝试点用小灰点；右上角小图给出累计虚拟时间随动作序号的阶梯曲线，并按动作类型着色。
  另可并排放一张对照图：$25$ 站旧版本在同一案例上的轨迹，直观展示路程差异。}}
  \caption{第三问典型局的行走轨迹与时间累积}\label{fig:q3-track}
\end{figure}

\cref{tab:q3paired}中收益最大的一行是测站布局（$47.04\%$），这与\cref{sec:analysis:fixed}的固定成本分析相吻合。特别地，混合贪心一行的成本分解显示：候选版平均\emph{少走} $5946\mmet$（约 $1189\msec$），代价是平均\emph{多做} $17.3$ 次检测与 $18.3$ 次失败清除（约 $141\msec$）。这正体现了\cref{sec:q1}以来的一贯取舍——\textbf{用少量额外测量换取大量移动节省}，而不是机械地减少动作次数。

\begin{table}[!htbp]
  \caption{第三问被否决的改进方案（均完成完整配对实验，清除率均为 $100\%$）}\label{tab:q3reject}\centering
  \begin{tabularx}{\textwidth}{>{\raggedright\arraybackslash}p{4.4cm} c c L}
    \toprule[1.5pt]
    方案 & 对照 s/源 & 候选 s/源 & 否决原因 \\
    \midrule[1pt]
    两步前瞻调度 $J(a)=C(p,a)+\lambda\min_bC(a,b)$ & $263.920$ & $292.795$ & 总移动 $12583\to14452\mmet$，$\mathrm{P}_{95}$ 亦恶化；复杂度上升而性能下降 \\
    非凸可行域（保留 \texttt{no\_signal} 挖洞后的非凸单元） & $263.920$ & $264.943$ & 单元碎片化，平均光学格数 $100.1\to102.3$，区域缩小未转化为更短路径 \\
    自适应第二测点（$5\times4$ 离散集合） & $263.920$ & $265.518$ & 均值与最坏值均恶化，见\cref{sec:q2} \\
    主动搜索（按未知区域覆盖收益选站） & $251.70$ & $260.95$ & 提前减少搜索后已发现目标测向不足，光学成本上升 \\
    \bottomrule[1.5pt]
  \end{tabularx}
\end{table}

\subsection{理论下界与优化空间}\label{sec:q3:lb}

为判断“还能优化多少”，本文建立一个下界模型：任何满足覆盖保证的策略都至少要
\textbf{(i)} 走遍 $7$ 个测站（下界取该站集的最短哈密顿路，目标处理段用随机点巡回的渐近公式 $0.7124\sqrt{NA}$ 估计不可避免的额外行程）；
\textbf{(ii)} 在所有站上排除空频道（未占用频道数 $\times$ 站数 $\times6\msec$，已发现频道平均扫半数站）；
\textbf{(iii)} 对每个源至少一次第二测点往返与一次成功清除。结果见\cref{tab:q3lb}。

\begin{table}[!htbp]
  \caption{第三问的时间下界与现有方案的差距（单位：虚拟秒）}\label{tab:q3lb}\centering
  \begin{tabular}{cccccccc}
    \toprule[1.5pt]
    源数 $n$ & 路线下界 & 扫描下界 & 每源附加 & 总下界 & 下界 s/源 & 现有本地 s/源 & 差距 \\
    \midrule[1pt]
    $10$ & $1874$ & $750$ & $310$ & $2934$ & $293$ & $310$ & $6\%$ \\
    $13$ & $2030$ & $700$ & $403$ & $3133$ & $241$ & $265$ & $9\%$ \\
    $16$ & $2180$ & $696$ & $496$ & $3372$ & $211$ & $226$ & $7\%$ \\
    \bottomrule[1.5pt]
  \end{tabular}
\end{table}

现有方案距下界只有 $6\%\sim9\%$，剩余差距来自巡站与目标处理交错的调度损失以及最坏情形余量。结合\cref{thm:seven}（站数不能再减）与\cref{cor:ring}（环半径已贴下界），本文认为\textbf{第三问的优化重点已从“搜索覆盖”转向“定位精度与光学覆盖开销”}，继续堆叠更复杂的调度器不再是有效方向。
